% Series RLC Circuit for Chapter 3 - Physical Laws
\begin{figure}[htbp]
\centering
\begin{circuitikz}[scale=1.0]
    % Circuit elements
    \draw (0,0)
        to[V, v=$v_{\text{in}}(t)$, invert] (0,3)
        to[R, l=$R$, v=$v_R$] (3,3)
        to[L, l=$L$, v=$v_L$] (6,3)
        to[C, l=$C$, v=$v_C$] (6,0)
        -- (0,0);

    % Current arrow
    \draw[->, thick, UofSCGarnet] (1.5,3.4) -- (2.5,3.4) node[midway, above] {$i(t)$};

    % Ground symbol
    \draw (3,0) node[ground] {};

    % Voltage annotations
    \node[right, font=\footnotesize] at (6.3,1.5) {Output: $v_C(t)$};
\end{circuitikz}
\caption{Series RLC circuit with input voltage $v_{\text{in}}(t)$ and output voltage $v_C(t)$ across the capacitor. Applying KVL: $v_{\text{in}} = v_R + v_L + v_C$.}
\label{fig:rlc_circuit}
\end{figure}
